The field of robotics has witnessed a significant paradigm shift from single-agent missions to complex Multi-Robot Systems (MRS) and Heterogeneous Multi-Agent Systems (MAS). 
Research in this area, which dates back to the late 1980s \cite{cao1997cooperative}, has demonstrated that coordinated teams of robots can achieve mission objectives that are beyond the capabilities of individual units. 
Heterogeneous systems, characterized by agents with diverse physical structures and capabilities, offer enhanced fault tolerance, scalability, and the ability to perform multi-faceted tasks through the fusion of multi-modal data.\newline\newline
The integration of Unmanned Aerial Vehicles (UAVs) and Unmanned Ground Vehicles (UGVs) is a prominent example of such synergy. While UAVs provide rapid exploration and global situational awareness from an aerial perspective, UGVs contribute with extended endurance, higher payload capacities, and stable interaction with the ground environment. 
This collaborative approach has found critical applications in structural monitoring of historical infrastructure, as well as in Search and Rescue (SAR)~\cite{HUANG2024109491}, industrial inspection~\cite{s20216384} and precision agriculture tasks.\newline\newline
Despite its potential, UAV-UGV collaboration faces significant scientific challenges~\cite{munasinghe2024uavugv}, particularly in optimizing communication and coordination under unstable network conditions. Recent research has explored solutions like the Unified Robotic Protocol (URP) to manage communication between diverse frameworks \cite{shahar2025ugv_uav}, and resource-efficient strategies like Semantic Communication (SemCom) \cite{zhang2023semantic}, which prioritizes the transmission of task-relevant semantic information over raw data to reduce bandwidth overhead.\newline\newline
The complexity of these systems necessitates high-fidelity simulation environments for safe and cost-effective validation. While traditional tools like Gazebo and Webots provide robust physics-based simulation, they often lack the photorealistic rendering required for advanced vision-based perception algorithms. In this context, simulators built on high-performance game engines, such as Microsoft's AirSim \cite{airsim}, have become state-of-the-art by offering realistic environmental physics and sensor behavior, as highlighted in recent comparative reviews \cite{nikolaiev2024comparative}.\newline\newline
Building upon this foundation, Cosys-AirSim~\cite{jansen2023cosys} was developed to address specific industrial and research needs, including more detailed sensor models and procedurally generated environments. However, a persistent limitation in existing frameworks, including the baseline AirSim implementation, is the lack of support for heterogeneous multi-vehicle configurations in unified scenes.
The proposed project addresses this gap by extending the Cosys-AirSim architecture to support concurrent interaction and control of UAVs and UGVs. This enhancement facilitates the integrated testing of perception and coordination strategies, positioning the platform as a leading tool for multi-domain autonomous research in immersive and virtual environments.
